% Subsection 3.3.1: 思维步骤标注

\texttt{thought\_step} 是 CoT 监督与 CoT 推理两条路径的共同信号，由 \texttt{annotate\_thought\_steps.py} 调用教师模型 \texttt{deepseek-chat} 离线生成。教师模型同时看到题目与参考解答，仅被要求产出「从题目自然导向该解答」的推理，避免把教师自身的解题误差传导到学生模型的监督信号中。

输出格式固定为 3$\sim$5 条带封闭标签的编号步骤：\texttt{Step N: [Tag] - explanation}，标签取自 \{Understand, Approach, Edge Cases, Implement, Verify\}，强制覆盖「理解题意—拟定算法—处理边界—实现—自检」关键子集。图~\ref{fig:ch3-thought-example} 展示与图~\ref{fig:ch3-sample} 同一题目的完整标注示例。

\begin{figure}[!htb]
\begin{footnotesize}
\begin{verbatim}
Step 1: [Understand]  - The goal is to partition the array into exactly k
        subsets where each subset's sum is unique. However, the provided
        solution actually checks for partitioning into k subsets with
        equal sums (total_sum / k), suggesting a possible misinterpretation.

Step 2: [Approach]    - For equal-sum partition, first check divisibility
        of total_sum by k; then recursively form subsets summing to
        target_sum = total_sum / k via backtracking over used elements.

Step 3: [Edge Cases]  - If total_sum % k != 0, return False immediately.
        Elements exceeding target_sum are pruned inside the backtracking.

Step 4: [Implement]   - subsets_with_sum() backtracks over unused elements
        to hit target_sum; can_partition() invokes it k times while
        sharing a single `used` array across the k subsets.

Step 5: [Verify]      - After k successful calls every element is used
        exactly once and every subset sums to target_sum (equal sums);
        note this does not guarantee the UNIQUE sums asked in the prompt.
\end{verbatim}
\end{footnotesize}
\vspace{-2mm}
\caption{一条 \texttt{thought\_step} 的完整结构（对应图~\ref{fig:ch3-sample} 中的样本）}
\label{fig:ch3-thought-example}
\end{figure}
